\documentclass[a4paper,12pt]{article}

\usepackage[T1]{fontenc}
\usepackage[italian]{babel}
\usepackage{graphicx}
\usepackage{geometry}
\usepackage{tabularx} % Per l'ambiente tabularx (tabelle)
\usepackage{calc} % Sempre per le tabelle
\usepackage[table]{xcolor}
\usepackage[most]{tcolorbox}
\usepackage{fancyhdr}
\usepackage[hidelinks, pdfnewwindow]{hyperref}
\usepackage{tikz}
\usepackage{pgf-pie}
\usepackage[utf8]{inputenc}
\usepackage{booktabs}
\usepackage{xcolor}  
\usepackage{float}
\usepackage[pdfnewwindow]{hyperref}
\usepackage{titlesec}
\usepackage{siunitx}
\usepackage{nicematrix}
\usepackage{colortbl}
\usepackage{eurosym}
\usepackage[dvipsnames]{xcolor}

\newcommand{\groupmail}{alittlebyte.swe@gmail.com}
\newcommand{\grouplogo}{../../../.github/site-src/assets/logo_alittlebyte.png}
\newcommand{\groupsymbol}{../../../.github/site-src/assets/solo_logo.png}

\geometry{
    left=2.5cm,
    right=2.5cm,
    top=2.5cm,
    bottom=2.5cm,
    headheight=331pt,
    includeheadfoot
}

\pagestyle{fancy}
\fancyhf{}

\renewcommand{\headrulewidth}{0pt}
\fancyhead{}

\renewcommand{\footrulewidth}{0.4pt}
\fancyfoot[L]{\includegraphics[height=0.6cm]{\groupsymbol}\hspace{0.45em}aLittleByte}
\fancyfoot[C]{\thepage}
\fancyfoot[R]{Verbale 2026-05-20}

\fancypagestyle{firstpage}{
    \fancyhf{}
    \renewcommand{\headrulewidth}{0pt}
    \renewcommand{\footrulewidth}{0.4pt}
    \fancyhead[C]{%
        \makebox[\textwidth][c]{%
            \begin{tabular}{c}
                \includegraphics[height=10cm]{\grouplogo}\\[1ex]
                \small\texttt{\groupmail}
            \end{tabular}%
        }
    }
    \fancyfoot[L]{\includegraphics[height=0.6cm]{\groupsymbol}\hspace{0.45em}aLittleByte}
    \fancyfoot[C]{\thepage}
    \fancyfoot[R]{Verbale 2026-05-20}
}

\providecommand{\glslink}[2]{\href{https://alittlebyte-19.github.io/Documentazione/glossario.html\#gls-#1}{\underline{#2}\textsuperscript{\scriptsize G}}}

\begin{document}

\thispagestyle{firstpage}

\vspace*{0.5cm}

\begin{center}
    {\LARGE \textbf{Verbale Esterno - 2026-05-20}}
\end{center}

\vspace{1.5cm}

\begin{center}
    \renewcommand{\arraystretch}{1.3}
    \begin{tcolorbox}[
        width=0.92\textwidth,
        colback=gray!6,
        colframe=gray!45,
        boxrule=0.5pt,
        arc=2mm,
        boxsep=0pt,
        left=8pt, right=8pt, top=8pt, bottom=8pt
    ]
        \begin{tabularx}{\linewidth}{>{\bfseries}p{3.5cm} X}
            Gruppo      & aLittleByte (gruppo 19) \\
            Redattore   & E. Bellet \\
            Verificatori& A. Oloni\\
            Destinatari & Prof. Tullio Vardanega, Prof. Riccardo Cardin \\
            Data        & 2026-05-20\\
        \end{tabularx}
    \end{tcolorbox}
\end{center}

\newpage

\newgeometry{left=2.5cm,right=2.5cm,top=2.5cm,bottom=2.5cm,includefoot}
\tableofcontents
\newpage

\section{Informazioni Generali}
\section*{Specifiche Documento}
\hrule
\vspace{1cm}

\begin{tabular}{>{\bfseries}p{5cm} | p{10cm}}
Luogo Riunione & Google Meet \\[6pt]

Orario (Inizio - Fine) & 14:30 -- 14.50 \\[6pt]

\glslink{redattore}{Redattore} &
E. Bellet\\[6pt]

\glslink{amministratore}{Amministratore} & 
D. Belluz\\[6pt]

\glslink{verificatore}{Verificatore} &
A. Oloni\\[6pt]

Partecipanti &
  A. Oloni\\ &
  A. Gastaldello\\ &
  E. Bellet\\ & 
  D. Belluz\\ &
  L. Soligo\\[6pt]

  Partecipanti Esterni &
  Luca Iuzzolino\\ &
  GianPaolo Ferrarin\\ [6pt]

\end{tabular}


\section{Ordine del Giorno}
    \begin{itemize}
    \item Mostrare le tecnologie usate
    \item Esposizione del PoC
    \item Conclusioni e Richiesta del Bucket
    \end{itemize}


\section{Attività svolte}
\subsection{Esposizione del PoC}
\textbf{Angelica Gastaldello}, con la presenza della maggior parte del gruppo, ha mostrato il \textit{Proof of Concept} all'azienda, nel quale sono stati mostrati i principali casi d'uso, il quale l'azienda ha dato esito positivo, sia a livello di utilizzo che per quanto riguarda l'interfaccia grafica.

\subsection{Spiegazione delle tecnologie usate}
\textbf{Leonardo Soligo} nell'esposizione ha spiegato le varie tecnologie usate, ossia:
\begin{itemize}
    \item Lo stack tecnologico include \glslink{redis}{Redis} per \glslink{metadati}{metadati} e cache, e \glslink{minio}{Minio} per lo storage di documenti originali e suddivisi.
    \item È stata sviluppata un'interfaccia admin per la gestione delle credenziali e per simulare risposte \glslink{intelligenza-artificiale-ai}{AI}, testando \glslink{workflow}{workflow} di keytab.
    \item L'elaborazione dei PDF viene delegata a un worker tramite coda \glslink{redis}{Redis}, senza bloccare la richiesta, e i documenti vengono salvati su \glslink{minio}{Minio}.
    \item L'assistente generativo permette di impostare tono e stile per generare bozze, con l'elaborazione dei PDF che delega il contenuto testuale a Bedrock.
    \item Sono emerse limitazioni nell'estrazione dati dai PDF, potenzialmente legate all'uso di Bedrock rispetto a un \glslink{ocr-riconoscimento-ottico-dei-caratteri}{OCR} più preciso, e si è discusso della necessità di identificare l'utente destinatario.
\end{itemize}

\subsection{Conclusioni e Richiesta del Bucket}
La dimostrazione ha avuto successo nel mostrare le funzionalità dell'applicazione e l'integrazione delle tecnologie. È stata confermata la disponibilità a fornire il bucket S3 richiesto per facilitare i test futuri. Ulteriori sviluppi potrebbero includere l'identificazione dell'utente e l'affinamento dei \glslink{prompt}{prompt} per migliorare l'estrazione dei dati dai documenti.

\section{To-Do List}
\begin{table}[h]
    \centering
    \begin{tabular}{|p{4.5cm}|p{5cm}|l|}
    \hline
        \rowcolor{lightgray}\textit{\textbf{Persona Incaricata}}  & \textbf{Task Assegnata} &\textbf{Link \glslink{issue}{Issue}} \\
    \hline
        \textit{Eleonora Bellet} & Redigere il verbale 2026-05-20 & \href{https://github.com/aLittleByte-19/Documentazione/issues/94}{\#94} \\
    \hline
    \end{tabular}
    \label{tab:placeholder}
\end{table}

\end{document}
